\chapter{Approximation and Online Algorithms}\label{ch:approx}

\begin{center}\itshape
When the best answer is out of reach, prove how far you are from it;\\
when the future is unknown, prove how much it can cost you.\\
\upshape(the two compromises of this chapter, in our words)
\end{center}

In 1976 two mathematicians on opposite sides of the iron curtain found the same algorithm for the travelling salesman problem. Nicos Christofides, at Imperial College London, described it in a report of the Graduate School of Industrial Administration of Carnegie Mellon University: a tour through $n$ cities, with distances satisfying the triangle inequality, never more than one and a half times as long as the shortest one.
Anatoliy Serdyukov, in Novosibirsk, found it the same year and published it in 1978 in a Russian journal, \emph{Upravlyaemye Sistemy}, where it remained unknown to most of the West. In 2020 René van Bevern and Viktoriia Slugina told the story in \emph{Historia Mathematica}, and many authors now speak of the Christofides--Serdyukov algorithm. For forty-five years nobody could beat the factor $\frac32$ in general; in 2020 Anna Karlin, Nathan Klein and Shayan Oveis Gharan finally did, by $10^{-36}$.

The travelling salesman problem is NP-hard: Tome~C3 will show that a polynomial algorithm for it would give one for every problem whose solutions can be checked quickly, which almost nobody believes possible. This chapter is about what one can still guarantee. An \emph{approximation algorithm} runs in polynomial time and returns a solution provably within a factor of the optimum, although the optimum itself is unknown; the art is to find a quantity that bounds the optimum and that the algorithm can compare itself to. We meet the main techniques: a combinatorial lower bound for vertex cover,
greedy choices paid by prices for set cover, trees and matchings for the travelling salesman, rounding for knapsack, linear programs and randomised rounding for satisfiability, and a semidefinite program with a random hyperplane for the maximum cut. Then we turn to \emph{online algorithms}, which must decide before they know the future: when to stop renting skis, which page to evict from a cache, and how to follow the advice of experts. The probability is that of Tome~3; linear programming is Chapter~\ref{ch:lp}.

\begin{think}
You must place fire stations at some of the crossroads of a city so that every street has a station at one of its ends, and you want few stations. Take any street with no station at either end, and build a station at \emph{both} of its ends; repeat. You may build too many. Can you prove that you build at most twice as many as necessary, without knowing how many are necessary?
\end{think}

\section{Approximation ratios and vertex cover}\label{sec:vc}

\begin{definition}
Let $\alpha\ge1$. An algorithm for a minimisation problem is an \emph{$\alpha$-approximation} if it runs in polynomial time and returns, on every instance, a feasible solution of cost at most $\alpha\cdot\mathrm{OPT}$. For a maximisation problem, with $\alpha\le1$, it returns a value at least $\alpha\cdot\mathrm{OPT}$. A randomised algorithm achieves the ratio \emph{in expectation} if the expected cost (or value) satisfies the inequality.
A \emph{polynomial-time approximation scheme} (PTAS) is a family of $(1+\varepsilon)$-approximations, one for each $\varepsilon>0$; it is \emph{fully polynomial} (FPTAS) if its running time is polynomial in the size of the input and in $1/\varepsilon$.
\end{definition}

A \emph{vertex cover} of a graph $G=(V,E)$ is a set of vertices meeting every edge; finding a smallest one is NP-hard (Tome~C3). The rule of the box is the following algorithm, attributed to Fănică Gavril and Mihalis Yannakakis in Garey and Johnson's book of 1979.

\begin{proposition}\label{prop:vcmatching}
Let $M$ be a maximal matching of $G$ (one to which no edge can be added). The set of the $2|M|$ endpoints of $M$ is a vertex cover of size at most $2\,\mathrm{OPT}$.
\end{proposition}
\begin{proof}
An edge with no endpoint in the set could be added to $M$, contradicting maximality, so the set is a cover. Any vertex cover contains an endpoint of each edge of $M$, and these endpoints are distinct since the edges of $M$ are disjoint: $\mathrm{OPT}\ge|M|$.
\end{proof}

The proof compared the algorithm to a lower bound on the optimum, the size of a matching, and never to the optimum itself. For vertices with weights $w_v\ge0$, a matching is no longer a good lower bound, but a linear program is:
\[
\mathrm{LP}=\min\Bigl\{\sum_vw_vx_v:\ x_u+x_v\ge1\text{ for }uv\in E,\ x\ge0\Bigr\}\le\mathrm{OPT},
\]
since the indicator vector of any cover is feasible.

\begin{proposition}\label{prop:vclp}
Let $x$ be an optimal solution of the linear program, computed in polynomial time (Section~\ref{sec:lpspeed}). Then $C=\{v:x_v\ge\frac12\}$ is a vertex cover of weight at most $2\,\mathrm{LP}\le2\,\mathrm{OPT}$.
\end{proposition}
\begin{proof}
For every edge $uv$, $x_u+x_v\ge1$, so one of them is at least $\frac12$: $C$ is a cover. And $\sum_{v\in C}w_v\le\sum_{v\in C}2w_vx_v\le2\,\mathrm{LP}$.
\end{proof}

Can one do better than $2$? Irit Dinur and Shmuel Safra proved in 2005 that a $1.36$-approximation would imply $\mathsf P=\mathsf{NP}$, and Subhash Khot and Oded Regev proved in 2008 that $2-\varepsilon$ is impossible for every $\varepsilon>0$ if Khot's \emph{unique games conjecture} holds. The humble factor $2$ may well be the truth.

\section{Set cover and the price of greed}\label{sec:setcover}

Let $U$ be a set of $n$ elements and $S_1,\dots,S_k$ subsets covering $U$, with costs $c_j>0$; we want a subfamily covering $U$ of least total cost. Vertex cover is the special case where the elements are edges and the sets are the stars of the vertices. The \emph{greedy algorithm} repeatedly chooses the set that covers new elements at the lowest cost per new element, that is which minimises $c_j/|S_j\setminus\text{covered}|$, until everything is covered. Let $H_d=1+\frac12+\dots+\frac1d$.

\begin{theorem}[Johnson, Lovász, Chvátal]\label{thm:setcover}
If every set has at most $d$ elements, the greedy algorithm returns a cover of cost at most $H_d\,\mathrm{OPT}\le(1+\ln n)\,\mathrm{OPT}$.
\end{theorem}
\begin{proof}
When greedy chooses a set $S$ that covers the new elements $N$, charge each $e\in N$ the \emph{price} $c(S)/|N|$. Each element is charged once, so the total cost of the greedy cover is the sum of all prices. We claim that for every set $S$ of the instance, $\sum_{e\in S}\mathrm{price}(e)\le H_{|S|}\,c(S)$.
Order the elements of $S$ as $e_1,\dots,e_s$ in the order in which greedy covers them (ties broken arbitrarily). Just before $e_i$ is covered, at least $s-i+1$ elements of $S$ are uncovered, so $S$ itself would cover them at cost per element at most $c(S)/(s-i+1)$; greedy chose a set at least as cheap per element, so $\mathrm{price}(e_i)\le c(S)/(s-i+1)$. Summing over $i$ gives $c(S)(\frac1s+\dots+\frac11)=H_sc(S)$.
Now let $\mathcal O$ be an optimal cover. Every element lies in some set of $\mathcal O$, so the total price is at most $\sum_{S\in\mathcal O}\sum_{e\in S}\mathrm{price}(e)\le H_d\sum_{S\in\mathcal O}c(S)=H_d\,\mathrm{OPT}$. Finally $H_d\le H_n\le1+\ln n$.
\end{proof}

The prices have a meaning. Divided by $H_d$, they form a feasible solution $y$ of the dual of the linear relaxation of set cover, $\max\{\sum_ey_e:\sum_{e\in S}y_e\le c(S)\text{ for all }S,\ y\ge0\}$, which is exactly what the claim says; by weak duality (Proposition~\ref{prop:weakduality}), the greedy cost is at most $H_d$ times the value of the linear program. This technique, called \emph{dual fitting}, proves more than the theorem: greedy is within $H_d$ of a lower bound that can be computed. It cannot be improved much. Uriel Feige showed in 1998 that a $(1-\varepsilon)\ln n$-approximation would put $\mathsf{NP}$ in quasi-polynomial time, and Dinur and David Steurer showed in 2014 that it would imply $\mathsf P=\mathsf{NP}$:
the greedy algorithm, one of the first approximation algorithms ever analysed (David Johnson, 1974), is essentially optimal.

\section{The travelling salesman}\label{sec:tsp}

An instance of the \emph{travelling salesman problem} (TSP) is a complete graph on $n$ vertices with nonnegative edge lengths $d(u,v)$; a \emph{tour} is a cycle through all the vertices, and we want a shortest one. Without assumptions on the lengths, nothing can be guaranteed.

\begin{proposition}\label{prop:tsphard}
Let $\alpha(n)\ge1$ be computable in time polynomial in $n$. If there is a polynomial $\alpha(n)$-approximation for the TSP, then one can decide in polynomial time whether a graph has a Hamiltonian cycle, and so $\mathsf P=\mathsf{NP}$.
\end{proposition}
\begin{proof}
Given a graph $G$ on $n$ vertices, give length $1$ to its edges and $\lceil\alpha(n)\,n\rceil+1$ to the other pairs. If $G$ has a Hamiltonian cycle, the shortest tour has length $n$, and the approximation returns a tour of length at most $\alpha(n)n$. If not, every tour uses a non-edge and has length more than $\alpha(n)n$. So comparing the returned length with $\alpha(n)n$ decides Hamiltonicity, a problem that Karp showed NP-complete in 1972 (Tome~C3).
\end{proof}

So we assume the \emph{triangle inequality} $d(u,w)\le d(u,v)+d(v,w)$, as for distances on a map; this is the \emph{metric} TSP.

\begin{lemma}[shortcutting]\label{lem:shortcut}
In a metric instance, every closed walk that visits all vertices can be turned into a tour that is not longer, by visiting the vertices in the order of their first appearance in the walk.
\end{lemma}
\begin{proof}
Between two consecutive first appearances $u$ and $v$ the walk follows a path from $u$ to $v$, which by the triangle inequality, applied repeatedly, is at least as long as the edge $uv$ of the tour; the same holds for the last vertex and the return to the start.
\end{proof}

A shortest tour minus one edge is a spanning path, hence a spanning tree; so the minimum spanning tree $T$ (Chapter~\ref{ch:graphs}) satisfies $w(T)\le\mathrm{OPT}$. Doubling every edge of $T$ gives a connected multigraph with all degrees even, hence a closed walk through every edge (Euler's theorem, Tome~C1, Theorem~5.4), of length $2w(T)$, and shortcutting it gives a $2$-approximation. The Christofides--Serdyukov algorithm is smarter about how to make the degrees even (Figure~\ref{fig:christofides}).

\begin{figure}[ht]
\centering
\fig{ch12-christofides}
\caption{The Christofides--Serdyukov algorithm on nine points in the plane. Left: a minimum spanning tree, with its vertices of odd degree in red. Middle: a minimum perfect matching of the odd vertices (dashed) makes all degrees even. Right: the tour obtained by shortcutting an Euler tour of tree plus matching (green), and a shortest tour (grey) for comparison.}
\label{fig:christofides}
\end{figure}

\begin{theorem}[Christofides, Serdyukov, 1976]\label{thm:christofides}
For the metric TSP, the following is a $\frac32$-approximation. Compute a minimum spanning tree $T$; let $O$ be the set of vertices of odd degree in $T$; compute a minimum-weight perfect matching $M$ on $O$; find an Euler tour of the multigraph $T\cup M$ and shortcut it.
\end{theorem}
\begin{proof}
By the handshake lemma $|O|$ is even, so $M$ exists; it can be computed in polynomial time by Edmonds' weighted matching algorithm (1965), which we only quote. In $T\cup M$ every vertex of $O$ gains one edge and the others none, so all degrees are even, and the multigraph is connected: it has an Euler tour of length $w(T)+w(M)$, and the final tour is not longer (Lemma~\ref{lem:shortcut}).
We know $w(T)\le\mathrm{OPT}$. Shortcut an optimal tour to visit only the vertices of $O$, in the same cyclic order: this gives a cycle $C$ on $O$ of length at most $\mathrm{OPT}$. As $|O|$ is even, the edges of $C$ taken alternately form two perfect matchings of $O$, of total length $w(C)$, so one of them has length at most $\frac12\mathrm{OPT}$, and $w(M)\le\frac12\mathrm{OPT}$. Altogether the tour has length at most $\frac32\mathrm{OPT}$.
\end{proof}

The analysis is tight for some instances, and the algorithm of Karlin, Klein and Oveis Gharan replaces the minimum spanning tree by a random spanning tree drawn from a carefully chosen distribution to save its $10^{-36}$. For points in the Euclidean plane, Sanjeev Arora and Joseph Mitchell found independently in 1996 a polynomial-time approximation scheme: there, any ratio $1+\varepsilon$ is possible.

\section{Knapsack: an approximation scheme}\label{sec:fptas}

The knapsack problem of Chapter~\ref{ch:dynamic} (items with weights $w_i$ and values $v_i$, a capacity $W$; maximise the value of a subset of weight at most $W$) is NP-hard too, but in a weak sense: Exercise~\ref{ex:knapsackvalue} solves it exactly in time $O(n^2v_{\max})$, which is polynomial when the values are small. The idea of Oscar Ibarra and Chul Kim (1975) is to make the values small by rounding them, and to control the damage.

\begin{theorem}[Ibarra--Kim]\label{thm:fptas}
Assume every item fits alone ($w_i\le W$). For $\varepsilon>0$, let $K=\varepsilon v_{\max}/n$ and $v'_i=\lfloor v_i/K\rfloor$, and solve the knapsack problem exactly for the values $v'_i$. The subset found has value at least $(1-\varepsilon)\mathrm{OPT}$, and the time is $O(n^3/\varepsilon)$. This is an FPTAS.
\end{theorem}
\begin{proof}
The rounded values are at most $v_{\max}/K=n/\varepsilon$, so the exact algorithm takes $O(n^2\cdot n/\varepsilon)$ time. Let $S'$ be the subset it finds and $S^*$ an optimal subset for the true values. Since $Kv'_i\le v_i<K(v'_i+1)$,
\[
\sum_{i\in S'}v_i\ge K\sum_{i\in S'}v'_i\ge K\sum_{i\in S^*}v'_i\ge\sum_{i\in S^*}(v_i-K)\ge\mathrm{OPT}-nK=\mathrm{OPT}-\varepsilon v_{\max}\ge(1-\varepsilon)\mathrm{OPT},
\]
where the second inequality is the optimality of $S'$ for the rounded values, and the last one holds because the most valuable item alone is feasible, so $\mathrm{OPT}\ge v_{\max}$.
\end{proof}

Strongly NP-hard problems, such as the TSP of Proposition~\ref{prop:tsphard}, have no FPTAS unless $\mathsf P=\mathsf{NP}$; the distinction between numbers written in binary and in unary is the subject of Tome~C3.

\section{Satisfiability: randomised rounding}\label{sec:maxsat}

In MAX-SAT we are given clauses $C_1,\dots,C_m$, each a disjunction of literals ($x_i$ or $\bar x_i$) on distinct variables, with weights $w_C\ge0$, and we look for a truth assignment maximising the total weight of the satisfied clauses. In 1974 Johnson observed the simplest bound.

\begin{proposition}\label{prop:maxsatrandom}
A uniformly random assignment satisfies a clause with $k$ literals with probability $1-2^{-k}$. So its expected value is at least $\frac12\sum_Cw_C\ge\frac12\mathrm{OPT}$, and at least $\frac78\mathrm{OPT}$ if every clause has three literals.
\end{proposition}
\begin{proof}
The clause fails only if its $k$ literals, which involve distinct variables, are all false, which has probability $2^{-k}$. Use linearity of expectation, and $\mathrm{OPT}\le\sum_Cw_C$.
\end{proof}

The method of conditional expectations (Proposition~\ref{prop:condexp}) turns this into a deterministic algorithm (Exercise~\ref{ex:maxsatderand}). Johan Håstad proved in 2001 that for clauses of exactly three literals, $\frac78+\varepsilon$ is NP-hard: the coin flip is optimal. For short clauses, the coin is weak (a clause $x_1$ alone is satisfied half the time), and a linear program helps. With $C^+$ and $C^-$ the sets of variables appearing positively and negatively in $C$, consider
\[
\mathrm{LP}=\max\Bigl\{\sum_Cw_Cz_C:\ z_C\le\sum_{i\in C^+}y_i+\sum_{i\in C^-}(1-y_i),\ 0\le y_i\le1,\ 0\le z_C\le1\Bigr\}\ge\mathrm{OPT},
\]
which is a relaxation since an assignment gives a feasible point with $z_C=1$ exactly for the satisfied clauses.

\begin{lemma}\label{lem:concave}
For $k\ge1$ and $z\in[0,1]$, $1-(1-z/k)^k\ge\bigl(1-(1-1/k)^k\bigr)z$, and $1-(1-1/k)^k\ge1-\frac1e$.
\end{lemma}
\begin{proof}
The function $f(z)=1-(1-z/k)^k$ is concave on $[0,1]$ (its second derivative is $-\frac{k-1}k(1-z/k)^{k-2}\le0$), with $f(0)=0$; so on $[0,1]$ it lies above its chord, $f(z)\ge zf(1)$. And $(1-1/k)^k\le e^{-1}$ since $1-t\le e^{-t}$.
\end{proof}

\begin{theorem}[Goemans--Williamson, 1994]\label{thm:maxsat}
Let $(y,z)$ be an optimal solution of the linear program, and set each variable $x_i$ true independently with probability $y_i$. A clause $C$ with $k$ literals is then satisfied with probability at least $(1-(1-1/k)^k)z_C\ge(1-\frac1e)z_C$. Running this and the uniform random assignment, and keeping the better result, gives a $\frac34$-approximation in expectation.
\end{theorem}
\begin{proof}
The clause fails with probability $\prod_{i\in C^+}(1-y_i)\prod_{i\in C^-}y_i$, a product of $k$ numbers in $[0,1]$. By the inequality of arithmetic and geometric means it is at most
\[
\Bigl(\frac1k\Bigl(\sum_{i\in C^+}(1-y_i)+\sum_{i\in C^-}y_i\Bigr)\Bigr)^k=\Bigl(1-\frac1k\Bigl(\sum_{i\in C^+}y_i+\sum_{i\in C^-}(1-y_i)\Bigr)\Bigr)^k\le\Bigl(1-\frac{z_C}k\Bigr)^k,
\]
and Lemma~\ref{lem:concave} gives the first claim. For the second, the better of the two assignments is at least their average, whose expectation is at least $\sum_Cw_Cz_C\cdot\frac12\bigl[(1-2^{-k_C})+(1-(1-1/k_C)^{k_C})\bigr]$, using $1-2^{-k}\ge(1-2^{-k})z_C$. The bracket is $\frac12+1=\frac32$ for $k=1$, $\frac34+\frac34=\frac32$ for $k=2$, and at least $\frac78+1-\frac1e>1.5$ for $k\ge3$. So the expected value is at least $\frac34\mathrm{LP}\ge\frac34\mathrm{OPT}$.
\end{proof}

The constant $\frac34$ is the best possible for this linear program: the four clauses $x\vee y$, $x\vee\bar y$, $\bar x\vee y$, $\bar x\vee\bar y$ have $\mathrm{LP}=4$ (all $y_i=\frac12$) but $\mathrm{OPT}=3$ (Exercise~\ref{ex:maxsatgap}). To go beyond, one needs a stronger relaxation, which is the subject of the next section.

\section{Maximum cut: a random hyperplane}\label{sec:gw}

A \emph{cut} of a graph with edge weights $w_{ij}\ge0$ is a partition of the vertices into two sides, of value the total weight of the edges between them; MAX-CUT asks for the largest. A random partition cuts each edge with probability $\frac12$, which is a $\frac12$-approximation (Chapter~\ref{ch:randomised}), and for twenty years nothing better was known. In 1994 Michel Goemans and David Williamson replaced the two sides, $\pm1$, by unit vectors.

Write a cut as signs $s_i\in\{\pm1\}$; its value is $\sum_{i<j}w_{ij}\frac{1-s_is_j}2$. The \emph{semidefinite relaxation} replaces each $s_i$ by a unit vector $v_i\in\mathbb R^n$:
\[
\mathrm{SDP}=\max\Bigl\{\sum_{i<j}w_{ij}\frac{1-v_i\cdot v_j}2:\ \|v_i\|=1\Bigr\}\ge\mathrm{OPT},
\]
since $v_i=s_ie_1$ is feasible. The matrix $X_{ij}=v_i\cdot v_j$ is a Gram matrix, that is a positive semidefinite matrix with ones on the diagonal, and every such matrix is a Gram matrix (take the columns of its square root, Tome~S1, Chapter~5). So the relaxation maximises a linear function of $X$ over a convex set of matrices; such \emph{semidefinite programs} can be solved to within any $\varepsilon>0$ in time polynomial in the input and in $\log(1/\varepsilon)$, by the ellipsoid method or by interior-point methods, a result we quote. We ignore this $\varepsilon$ below.

\begin{figure}[ht]
\centering
\fig{ch12-gw}
\caption{In the plane spanned by $v_i$ and $v_j$, a random hyperplane through $0$ is a random line, with direction uniform. It separates $v_i$ from $v_j$ exactly when its normal $r$ points into one of the two shaded sectors of angle $\theta$: probability $2\theta/2\pi=\theta/\pi$.}
\label{fig:gw}
\end{figure}

\begin{lemma}[random hyperplane]\label{lem:hyperplane}
Let $r\in\mathbb R^n$ have independent standard normal coordinates, and let $v_i,v_j$ be unit vectors at angle $\theta\in[0,\pi]$. Then $\mathbb P\bigl(\operatorname{sign}(r\cdot v_i)\ne\operatorname{sign}(r\cdot v_j)\bigr)=\theta/\pi$.
\end{lemma}
\begin{proof}
If $v_j=\pm v_i$ the claim is clear ($r\cdot v_i=0$ has probability $0$). Otherwise let $u_1,u_2$ be an orthonormal basis of the plane $\Pi$ spanned by $v_i,v_j$, completed to an orthonormal basis $u_1,\dots,u_n$ of $\mathbb R^n$. The density $(2\pi)^{-n/2}e^{-\|x\|^2/2}$ of $r$ is unchanged by orthogonal maps, so the coordinates $(r\cdot u_1,\dots,r\cdot u_n)$ are again independent standard normals.
The signs depend only on the projection $\rho=(r\cdot u_1)u_1+(r\cdot u_2)u_2$ of $r$ on $\Pi$, since $r\cdot v=\rho\cdot v$ for $v\in\Pi$. The pair $(r\cdot u_1,r\cdot u_2)$ has the density $\frac1{2\pi}e^{-(a^2+b^2)/2}$, which depends only on the radius, so in polar coordinates the angle of $\rho$ is uniform on the circle. Finally $\rho\cdot v_i$ and $\rho\cdot v_j$ have different signs exactly when the direction of $\rho$ lies within $90^\circ$ of one of the two vectors and not of the other, which happens on two arcs of angle $\theta$ each (Figure~\ref{fig:gw}): probability $2\theta/2\pi$.
\end{proof}

\begin{lemma}\label{lem:gwconstant}
For all $\theta\in[0,\pi]$, $\dfrac\theta\pi\ge0.878\cdot\dfrac{1-\cos\theta}2$.
\end{lemma}
\begin{proof}
Let $\alpha=0.878$ and $h(\theta)=\frac\theta\pi-\frac\alpha2(1-\cos\theta)$, so $h(0)=0$ and $h'(\theta)=\frac1\pi-\frac\alpha2\sin\theta$. Let $s=\frac2{\alpha\pi}=0.725079\ldots<1$, $\theta_1=\arcsin s$ and $\theta_2=\pi-\theta_1$. Then $h'>0$ on $[0,\theta_1)$ and $(\theta_2,\pi]$, and $h'<0$ on $(\theta_1,\theta_2)$; so the minimum of $h$ on $[0,\pi]$ is $\min(h(0),h(\theta_2))$.
Now $\theta_1=0.811149\ldots$, $\theta_2=2.330442\ldots$, $\cos\theta_2=-\sqrt{1-s^2}=-0.688665\ldots$, and $h(\theta_2)=0.741803\ldots-0.741324\ldots=0.000479\ldots>0$. (Each value is given to six digits; the errors are far below the margin.) So $h\ge0$ on $[0,\pi]$.
\end{proof}

\begin{theorem}[Goemans--Williamson]\label{thm:gw}
Solve the semidefinite relaxation, draw $r$ as in Lemma~\ref{lem:hyperplane}, and put $i$ on the side of the sign of $r\cdot v_i$. The expected value of the cut is at least $0.878\,\mathrm{SDP}\ge0.878\,\mathrm{OPT}$.
\end{theorem}
\begin{proof}
With $\theta_{ij}$ the angle between $v_i$ and $v_j$, so that $v_i\cdot v_j=\cos\theta_{ij}$, linearity of expectation and the two lemmas give
$\mathbb E[\text{cut}]=\sum_{i<j}w_{ij}\frac{\theta_{ij}}\pi\ge0.878\sum_{i<j}w_{ij}\frac{1-\cos\theta_{ij}}2=0.878\,\mathrm{SDP}$.
\end{proof}

The best constant in Lemma~\ref{lem:gwconstant} is $\alpha_{\mathrm{GW}}=\min_\theta\frac{2\theta}{\pi(1-\cos\theta)}=0.87856\ldots$. Håstad showed that approximating MAX-CUT within $\frac{16}{17}\approx0.941$ is NP-hard, and in 2007 Khot, Guy Kindler, Elchanan Mossel and Ryan O'Donnell showed that if the unique games conjecture holds, $\alpha_{\mathrm{GW}}+\varepsilon$ is NP-hard: the random hyperplane would then be the last word. The pentagon shows the limit of the relaxation itself (Exercise~\ref{ex:pentagon}).

\section{Online algorithms}\label{sec:online}

An \emph{online} algorithm receives its input piece by piece and must act on each piece before seeing the next. It is compared with the optimal \emph{offline} solution, which knows the whole input in advance. An online algorithm for a minimisation problem is \emph{$c$-competitive} if there is a constant $b$ with $\mathrm{ALG}(\sigma)\le c\cdot\mathrm{OPT}(\sigma)+b$ for every input sequence $\sigma$. The idea goes back to Daniel Sleator and Robert Tarjan in 1985, and the name to Anna Karlin, Mark Manasse, Larry Rudolph and Sleator in 1988.

The simplest example is \emph{ski rental}. Skis cost $1$ per day to rent and $B$ to buy, where $B\ge1$ is an integer, and you do not know how many days the season will last. Every deterministic strategy is of the form ``rent until day $t-1$, buy on day $t$ if the season is still on'' (or never buy).

\begin{proposition}\label{prop:ski}
Buying on day $B$ costs at most $(2-\frac1B)\mathrm{OPT}$, and no deterministic strategy has a better ratio.
\end{proposition}
\begin{proof}
If the season lasts $D$ days, $\mathrm{OPT}=\min(D,B)$. With $t=B$: if $D<B$ we pay $D=\mathrm{OPT}$; if $D\ge B$ we pay $(B-1)+B=2B-1$ while $\mathrm{OPT}=B$. Conversely, against the strategy that buys on day $t$, let the season end on day $t$, just after the purchase: the strategy pays $t-1+B$ and $\mathrm{OPT}=\min(t,B)$. If $t\le B$ the ratio is $1+\frac{B-1}t\ge2-\frac1B$; if $t>B$ it is $\frac{t-1+B}B\ge2$. Never buying has unbounded ratio.
\end{proof}

A randomised strategy, buying on a random day with a well-chosen distribution, achieves the ratio $\frac e{e-1}\approx1.58$ against an adversary who does not see the coins, and this is optimal (Karlin, Manasse, McGeoch and Owicki, 1994). The same reasoning, wait until what you have lost equals what it would have cost to commit, governs when to spin down a disk, when to stop waiting for a lock, or when to take a taxi instead of waiting for the bus.

\subsection*{Paging}

A cache holds $k$ pages; requests for pages arrive one by one; a request for a page not in the cache is a \emph{fault}, and the requested page must be brought in, evicting another one if the cache is full. László Belady showed in 1966 that the offline optimum evicts the page whose next request is furthest in the future. Online, the classic rule is LRU: evict the \emph{least recently used} page.

\begin{theorem}[Sleator--Tarjan]\label{thm:paging}
\begin{enumerate}
\item For every request sequence, $\mathrm{LRU}\le k\cdot\mathrm{OPT}+k$, where both start with the same cache and $\mathrm{OPT}$ is the least number of faults of any algorithm.
\item For every deterministic online algorithm $A$ and every $N$, there is a sequence of $N$ requests to $k+1$ pages on which $A$ faults $N$ times and $\mathrm{OPT}$ at most $\lceil N/k\rceil$ times.
\end{enumerate}
\end{theorem}
\begin{proof}
(1) Cut the sequence into \emph{phases}: each phase is a maximal block of consecutive requests containing at most $k$ distinct pages, so that each new phase starts with a request for a page not seen in the previous phase, which then has exactly $k$ distinct pages. Within a phase, LRU faults at most once on each page: once requested in the phase, a page $p$ could be evicted only if it were the least recently used, which requires $k$ distinct other pages to have been requested after it, hence $k+1$ distinct pages in the phase. So LRU faults at most $k$ times per phase.
Now consider phase $i$ and the first request of phase $i+1$. Let $p$ be the first page requested in phase $i$. The later requests of phase $i$, together with the first request of phase $i+1$, involve $k$ distinct pages different from $p$: the $k-1$ other pages of phase $i$, and the new page that starts phase $i+1$. Just after serving the first request of phase $i$, any algorithm has $p$ in its cache, so at most $k-1$ of these $k$ pages; it faults at least once before the end of this stretch. These stretches are disjoint for different $i$, so $\mathrm{OPT}\ge(\text{number of phases})-1$, while $\mathrm{LRU}\le k\cdot(\text{number of phases})$.

(2) With $k+1$ pages, one of them is always missing from the cache of $A$; request it every time. Then $A$ faults on every request. An offline algorithm, on each fault, can evict a page that is not requested in the next $k-1$ requests (those involve at most $k-1$ pages, and the cache holds $k$ pages other than the requested one); then its next fault is at least $k$ requests later. So it faults at most $\lceil N/k\rceil$ times.
\end{proof}

The ratio $k$ looks terrible, and is far from what LRU does in practice, where requests have locality; the competitive analysis is a worst case over all sequences. Amos Fiat, Richard Karp, Michael Luby, Lyle McGeoch, Sleator and Neal Young showed in 1991 that a randomised \emph{marking} algorithm is $2H_k$-competitive, exponentially better, and that $H_k$ is a lower bound for randomised algorithms.

\subsection*{Learning from experts}

Every morning, $n$ experts give advice; you follow a mixture of them, and in the evening you learn the loss $\ell_t(i)\in[0,1]$ of each expert $i$ on day $t$. Your loss is $\sum_ip_t(i)\ell_t(i)$ if you followed expert $i$ with probability $p_t(i)$. After $T$ days, you would like to have lost not much more than the best expert in hindsight, whose identity you could not know in advance. The \emph{multiplicative weights} (or \emph{Hedge}) algorithm, from the weighted majority algorithm of Nick Littlestone and Manfred Warmuth (1994) and the Hedge of Yoav Freund and Robert Schapire (1997), starts with weights $W_1(i)=1$, follows $p_t=W_t/\sum_iW_t(i)$, and updates $W_{t+1}(i)=W_t(i)e^{-\eta\ell_t(i)}$ for a parameter $\eta>0$.

\begin{theorem}[multiplicative weights]\label{thm:mw}
For every sequence of losses in $[0,1]$ and every expert $i$, with $L_i=\sum_t\ell_t(i)$, the loss of Hedge satisfies
\[
\sum_{t=1}^T p_t\cdot\ell_t\le\frac{\eta L_i+\ln n}{1-e^{-\eta}}.
\]
With $\eta=\sqrt{\ln n/T}\le1$, the \emph{regret} $\sum_tp_t\cdot\ell_t-\min_iL_i$ is at most $2\sqrt{T\ln n}+\ln n$.
\end{theorem}
\begin{proof}
Let $\Phi_t=\sum_iW_t(i)$, so $\Phi_1=n$. By convexity of $x\mapsto e^{-\eta x}$, $e^{-\eta x}\le1-(1-e^{-\eta})x$ for $x\in[0,1]$, so
\[
\Phi_{t+1}=\sum_iW_t(i)e^{-\eta\ell_t(i)}\le\Phi_t\bigl(1-(1-e^{-\eta})\,p_t\cdot\ell_t\bigr)\le\Phi_t\exp\bigl(-(1-e^{-\eta})\,p_t\cdot\ell_t\bigr).
\]
Hence $\Phi_{T+1}\le n\exp\bigl(-(1-e^{-\eta})\sum_tp_t\cdot\ell_t\bigr)$. On the other hand $\Phi_{T+1}\ge W_{T+1}(i)=e^{-\eta L_i}$. Taking logarithms, $(1-e^{-\eta})\sum_tp_t\cdot\ell_t\le\eta L_i+\ln n$.
For the regret, $1-e^{-\eta}\ge\eta-\frac{\eta^2}2=\eta(1-\frac\eta2)$ and $\frac1{1-\eta/2}\le1+\eta$ for $0<\eta\le1$, so $\sum_tp_t\cdot\ell_t\le(1+\eta)(L_i+\frac{\ln n}\eta)=L_i+\eta L_i+\frac{\ln n}\eta+\ln n$. With $L_i\le T$ and $\eta=\sqrt{\ln n/T}$ the extra terms are at most $2\sqrt{T\ln n}+\ln n$.
\end{proof}

The average regret per day tends to $0$ like $\sqrt{\ln n/T}$, whatever the losses, even if they are chosen by an adversary who knows the algorithm (but not its coins; a deterministic algorithm can be forced to have regret $T/2$, Exercise~\ref{ex:mwlower}). Few theorems have so many uses. Played against each other, two Hedge players converge to the value of a zero-sum game, which reproves von Neumann's minimax theorem (Exercise~\ref{ex:gamevalue}); with experts that are constraints, multiplicative weights solve linear programs approximately (Plotkin, Shmoys and Tardos, 1995); with experts that are training examples, they become the boosting algorithm AdaBoost of Freund and Schapire. Sanjeev Arora, Elad Hazan and Satyen Kale called it, in 2012, a meta-algorithm. Tome~C5 returns to it as online learning.

\begin{beauty}
An approximation guarantee is a proof about a number nobody knows. We bound the optimum from below by a matching, by a linear program, by a tree, by a cycle cut in two, by a semidefinite program, and we show that our solution is close to the bound; the optimum, squeezed in between, never needs to be computed. The online theorems do the same with time: the optimal hindsight, which no algorithm can have, becomes the yardstick, and the weights of the experts, multiplied a little each day, catch up with the best of them whatever the future holds.
\end{beauty}

\section{Exercises}

\begin{exercise}[$\star$]\label{ex:vcmatching}
Show that the factor $2$ in Proposition~\ref{prop:vcmatching} is attained on the complete bipartite graph $K_{n,n}$, and that the LP bound of Proposition~\ref{prop:vclp} cannot prove a ratio better than $2-\frac2n$ on the complete graph $K_n$.
\end{exercise}

\begin{exercise}[$\star\star$]\label{ex:vclphalf}
Show that every vertex of the polyhedron $\{x:x_u+x_v\ge1\ (uv\in E),\ x\ge0\}$ has all its coordinates in $\{0,\frac12,1\}$ (Nemhauser and Trotter).
\end{exercise}

\begin{exercise}[$\star\star$]\label{ex:setcovertight}
Let $k\ge2$ and $U$ be two rows of $2^k-1$ elements each, split into consecutive blocks of sizes $1,2,4,\dots,2^{k-1}$ in each row. The sets are the two rows, and for each $j$ the union $B_j$ of the two blocks of size $2^{j}$; all costs are $1$. Show that greedy may choose $B_{k-1},B_{k-2},\dots,B_0$, using $k$ sets instead of $2$.
\end{exercise}

\begin{exercise}[$\star\star$]\label{ex:kcenter}
\emph{$k$-center.} Given $n$ points in a metric space and an integer $k$, we want $k$ centers minimising the largest distance from a point to its nearest center. Show that choosing an arbitrary first center, then repeatedly the point farthest from the centers chosen so far, is a $2$-approximation (Gonzalez, 1985).
\end{exercise}

\begin{exercise}[$\star$]\label{ex:knapsackgreedy}
For the knapsack problem with every item fitting alone, sort the items by decreasing $v_i/w_i$ and take them in this order while they fit, stopping at the first one that does not. Show that the better of this set and the most valuable single item has value at least $\frac12\mathrm{OPT}$.
\end{exercise}

\begin{exercise}[$\star$]\label{ex:maxsatderand}
Turn Proposition~\ref{prop:maxsatrandom} into a deterministic algorithm that satisfies clauses of total weight at least $\frac12\sum_Cw_C$, by setting the variables one at a time.
\end{exercise}

\begin{exercise}[$\star$]\label{ex:maxsatgap}
For the four clauses $x\vee y$, $x\vee\bar y$, $\bar x\vee y$, $\bar x\vee\bar y$ with weight $1$, show that $\mathrm{OPT}=3$ and $\mathrm{LP}=4$.
\end{exercise}

\begin{exercise}[$\star\star$]\label{ex:pentagon}
For the $5$-cycle with unit weights, show that $\mathrm{OPT}=4$ and that $\mathrm{SDP}\ge\frac52(1+\cos\frac\pi5)\approx4.52$, using five vectors in the plane. What does this say about the ratio between the optimum and the relaxation?
\end{exercise}

\begin{exercise}[$\star\star$]\label{ex:fifo}
Show that FIFO (evict the page that entered the cache first) also satisfies $\mathrm{FIFO}\le k\cdot\mathrm{OPT}+k$.
\end{exercise}

\begin{exercise}[$\star$]\label{ex:mwlower}
Show that for every deterministic algorithm that follows one of two experts each day, there are losses in $\{0,1\}$ for which it loses $T$ while the best expert loses at most $T/2$.
\end{exercise}

\section{Hints}

\begin{hint}{ex:vcmatching}
In $K_{n,n}$ a maximal matching may be perfect. In $K_n$, consider $x\equiv\frac12$.
\end{hint}
\begin{hint}{ex:vclphalf}
If $x$ has coordinates outside $\{0,\frac12,1\}$, perturb those in $(0,\frac12)$ by $+\varepsilon$ and those in $(\frac12,1)$ by $-\varepsilon$, and the other way round.
\end{hint}
\begin{hint}{ex:setcovertight}
When $B_{k-1},\dots,B_{j+1}$ have been chosen, compare the number of new elements of $B_j$ and of each row.
\end{hint}
\begin{hint}{ex:kcenter}
Look at the $k$ centers and the point that realises the final radius: $k+1$ points, pairwise far apart.
\end{hint}
\begin{hint}{ex:knapsackgreedy}
The greedy set plus the first item that does not fit is worth at least the fractional optimum.
\end{hint}
\begin{hint}{ex:maxsatderand}
The conditional expectation, given the values of some variables, is a sum over clauses that can be computed exactly.
\end{hint}
\begin{hint}{ex:maxsatgap}
Every assignment falsifies exactly one clause.
\end{hint}
\begin{hint}{ex:pentagon}
Put $v_k$ at angle $\frac{4\pi k}5$.
\end{hint}
\begin{hint}{ex:fifo}
Within a phase, can FIFO fault twice on the same page?
\end{hint}
\begin{hint}{ex:mwlower}
Let the adversary punish the expert that the algorithm follows.
\end{hint}

\section{Solutions}

\begin{solution}{ex:vcmatching}
\idea Tightness examples show where the lower bound is weak.
\writeup In $K_{n,n}$ a perfect matching is maximal, and its $2n$ endpoints are all the vertices, while one side alone, of size $n$, is a cover: ratio $2$. In $K_n$ the point $x_v=\frac12$ for all $v$ is feasible, so $\mathrm{LP}\le\frac n2$, while every cover misses at most one vertex (two missing vertices would leave their edge uncovered), so $\mathrm{OPT}=n-1$. Any proof that compares with $\mathrm{LP}$ can therefore show no ratio better than $\frac{n-1}{n/2}=2-\frac2n$: the \emph{integrality gap} of this linear program tends to $2$.
\end{solution}

\begin{solution}{ex:vclphalf}
\idea A vertex is a point that is not the midpoint of two other points of the polyhedron.
\writeup Let $x$ be feasible with some coordinate outside $\{0,\frac12,1\}$; we write $x$ as the midpoint of two different feasible points. If some $x_v>1$, every constraint containing $x_v$ has slack and $x_v>0$, so $x\pm\varepsilon\mathbf 1_{\{v\}}$ are both feasible for small $\varepsilon>0$. Otherwise $0\le x\le1$; let $A=\{v:0<x_v<\frac12\}$ and $B=\{v:\frac12<x_v<1\}$, not both empty, and for small $\varepsilon>0$ let $x^\pm=x\pm\varepsilon(\mathbf 1_A-\mathbf 1_B)$. Nonnegativity holds for $\varepsilon$ small.
For an edge $uv$: if $x_u+x_v>1$, small changes keep the inequality. If $x_u+x_v=1$, then either both are $\frac12$ (unchanged), or one is $0$ and the other $1$ (unchanged), or one is in $(0,\frac12)$ and the other in $(\frac12,1)$, and then one gains $\pm\varepsilon$ and the other loses it, so the sum stays $1$. (A coordinate in $(0,\frac12)$ with an edge of sum exactly $1$ must have its neighbour in $(\frac12,1)$.) So $x=\frac12(x^++x^-)$ with $x^+\ne x^-$, and $x$ is not a vertex.
\end{solution}

\begin{solution}{ex:setcovertight}
\idea Each block set always looks slightly better than a row, because a row has already lost the blocks greedy took.
\writeup Each row has $1+2+\dots+2^{k-1}=2^k-1$ elements and $B_j$ has $2^{j+1}$. Suppose greedy has chosen $B_{k-1},\dots,B_{j+1}$ (none at the start, $j=k-1$). Then each row has $2^{j+1}-1$ uncovered elements, $B_j$ has $2^{j+1}$, and the other $B_i$ ($i<j$) have fewer. So $B_j$ is the best choice (strictly). After $B_0$ everything is covered: greedy uses $k$ sets, while the two rows cover $U$. Here $n=2(2^k-1)$, so the ratio $k/2$ is about $\frac12\log_2n$.
\end{solution}

\begin{solution}{ex:kcenter}
\idea Farthest-first produces points that are far from each other, and two of them must share an optimal cluster.
\writeup Let $c_1,\dots,c_k$ be the centers chosen, $r$ the final radius, and $p$ a point at distance $r$ from all the centers. When $c_j$ was chosen ($j\ge2$), it was the farthest point from $c_1,\dots,c_{j-1}$, at distance at least that of $p$ from them, which is at least $r$ (distances to a growing set of centers can only decrease). So the $k+1$ points $c_1,\dots,c_k,p$ are pairwise at distance at least $r$.
In an optimal solution with radius $\mathrm{OPT}$, two of these $k+1$ points are served by the same center $z$ (pigeonhole), and by the triangle inequality their distance is at most $2\,\mathrm{OPT}$. So $r\le2\,\mathrm{OPT}$.
\end{solution}

\begin{solution}{ex:knapsackgreedy}
\idea Greedy by density is optimal for the fractional problem, and it misses only the one item it had to cut.
\writeup Let the items be sorted, and let $j$ be the first that does not fit after $1,\dots,j-1$. The fractional optimum, in which items may be taken in part, takes $1,\dots,j-1$, a fraction of $j$, and nothing else (an exchange argument: replacing weight of a less dense item by weight of a denser one never decreases the value). So $\mathrm{OPT}\le\mathrm{OPT}_{\mathrm{frac}}\le v_1+\dots+v_{j-1}+v_j$.
Hence the larger of $v_1+\dots+v_{j-1}$ and $v_j\le v_{\max}$ is at least $\frac12\mathrm{OPT}$, and so is the larger of the greedy set and the best single item. Without the single item, greedy alone can be arbitrarily bad: one item of value $2$ and weight $1$, one of value $W$ and weight $W$.
\end{solution}

\begin{solution}{ex:maxsatderand}
\idea Never let the expected value drop.
\writeup Let $E(a_1,\dots,a_j)$ be the expected weight of satisfied clauses when $x_1,\dots,x_j$ are fixed to $a_1,\dots,a_j$ and the others are uniform. A clause already satisfied contributes $w_C$; a clause with all its literals fixed and false contributes $0$; otherwise it contributes $w_C(1-2^{-u})$ with $u$ its number of unfixed literals. So $E$ is computable in linear time. Since $E(a_1,\dots,a_j)$ is the average of $E(a_1,\dots,a_j,\mathrm{true})$ and $E(a_1,\dots,a_j,\mathrm{false})$, one of them is at least as large; choose it. Starting from $E()\ge\frac12\sum_Cw_C$, the final assignment has value $E(a_1,\dots,a_n)\ge\frac12\sum_Cw_C$.
\end{solution}

\begin{solution}{ex:maxsatgap}
\idea The four clauses forbid the four assignments of $(x,y)$, one each.
\writeup The clause $x\vee y$ is false only for $x=y=\mathrm{false}$, and similarly each clause is false for exactly one of the four assignments of $(x,y)$, a different one for each clause. So every assignment satisfies exactly three clauses: $\mathrm{OPT}=3$. With $y_x=y_y=\frac12$, each constraint reads $z_C\le\frac12+\frac12=1$, so $z_C=1$ for all four clauses and $\mathrm{LP}=4$ (the maximum possible). The ratio $\frac34$ shows that no rounding of this linear program can guarantee more than $\frac34\mathrm{LP}$.
\end{solution}

\begin{solution}{ex:pentagon}
\idea The vectors can do what signs cannot: go round an odd cycle with every step nearly opposite.
\writeup A cut of the $5$-cycle cannot cut all five edges, as the sides would alternate around an odd cycle; cutting four is possible (sides $+,-,+,-,-$). So $\mathrm{OPT}=4$. With $v_k=(\cos\frac{4\pi k}5,\sin\frac{4\pi k}5)$, consecutive vectors are at angle $\frac{4\pi}5$, so each edge contributes $\frac12(1-\cos\frac{4\pi}5)=\frac12(1+\cos\frac\pi5)$, and the total is $\frac52(1+\cos\frac\pi5)=\frac52\cdot\frac{5+\sqrt5}4\approx4.523$.
So $\mathrm{OPT}/\mathrm{SDP}\le4/4.523\approx0.884$: no rounding of this relaxation can guarantee more than $0.884$ times its value, not far from the $0.878$ of Goemans and Williamson. (Feige and Schechtman showed in 2002 that the gap of this relaxation is exactly $\alpha_{\mathrm{GW}}$, with other graphs.)
\end{solution}

\begin{solution}{ex:fifo}
\idea Same phases as for LRU; FIFO, too, never faults twice on a page within a phase.
\writeup Use the phases of the proof of Theorem~\ref{thm:paging}; the lower bound on $\mathrm{OPT}$ is unchanged. Suppose FIFO faults twice on a page $p$ within a phase. After the first fault $p$ entered the cache, and before the second it was evicted, so $k$ pages entered after $p$ in between (FIFO evicts the oldest entry, and the cache holds $k$ pages, so $p$ is evicted exactly at the $k$th insertion after its own). These $k$ insertions are of distinct pages different from $p$: a page $q$ inserted twice after $p$ would have been evicted in between, before $p$, which FIFO never does to a younger page.
All these requests lie in the phase, which would then contain $k+1$ distinct pages: a contradiction. So FIFO faults at most $k$ times per phase, and $\mathrm{FIFO}\le k\cdot\mathrm{OPT}+k$.
\end{solution}

\begin{solution}{ex:mwlower}
\idea A deterministic choice can be predicted, hence punished.
\writeup Each day the adversary computes, from the past, which expert the algorithm will follow, gives that expert loss $1$ and the other loss $0$. The algorithm loses $1$ every day, $T$ in total. The two experts' losses add up to $T$, so the better one has lost at most $T/2$: the regret is at least $T/2$, linear in $T$. Hedge escapes because the adversary, not knowing its coins, cannot make both experts bad at once.
\end{solution}

\section*{Sources and further reading}

\begin{itemize}
\item \emph{A research paper to read.} M.~X.~Goemans and D.~P.~Williamson, \emph{Improved approximation algorithms for maximum cut and satisfiability problems using semidefinite programming}, Journal of the ACM 42 (1995), 1115--1145 (conference version 1994). Read Section~2, which is Theorem~\ref{thm:gw} with the same two lemmas, and the discussion of how the semidefinite program is solved; then the section on MAX-2SAT, where the same rounding gives $0.878$ for clauses of two literals by representing \emph{true} as a special vector $v_0$. The paper is short, and it opened the field of semidefinite programming in approximation.
\item \emph{The stories.} N.~Christofides, \emph{Worst-case analysis of a new heuristic for the travelling salesman problem}, Report 388, Graduate School of Industrial Administration, Carnegie Mellon University, 1976; A.~I.~Serdyukov, \emph{On some extremal walks in graphs}, Upravlyaemye Sistemy 17 (1978) (in Russian); R.~van Bevern and V.~A.~Slugina, \emph{A historical note on the 3/2-approximation algorithm for the metric traveling salesman problem}, Historia Mathematica 53 (2020); A.~R.~Karlin, N.~Klein and S.~Oveis Gharan, \emph{A (slightly) improved approximation algorithm for metric TSP}, Proceedings of the 53rd ACM STOC (2021), 32--45.
\item D.~S.~Johnson, \emph{Approximation algorithms for combinatorial problems}, Journal of Computer and System Sciences 9 (1974), 256--278; V.~Chv\'atal, \emph{A greedy heuristic for the set-covering problem}, Mathematics of Operations Research 4 (1979), 233--235; U.~Feige, \emph{A threshold of $\ln n$ for approximating set cover}, Journal of the ACM 45 (1998), 634--652; O.~H.~Ibarra and C.~E.~Kim, \emph{Fast approximation algorithms for the knapsack and sum of subset problems}, Journal of the ACM 22 (1975), 463--468;
I.~Dinur and S.~Safra, \emph{On the hardness of approximating minimum vertex cover}, Annals of Mathematics 162 (2005), 439--485; S.~Khot and O.~Regev, \emph{Vertex cover might be hard to approximate to within $2-\varepsilon$}, Journal of Computer and System Sciences 74 (2008), 335--349; J.~H\aa stad, \emph{Some optimal inapproximability results}, Journal of the ACM 48 (2001), 798--859; S.~Khot, G.~Kindler, E.~Mossel and R.~O'Donnell, \emph{Optimal inapproximability results for MAX-CUT and other 2-variable CSPs?}, SIAM Journal on Computing 37 (2007), 319--357; T.~F.~Gonzalez, \emph{Clustering to minimize the maximum intercluster distance}, Theoretical Computer Science 38 (1985), 293--306.
\item L.~A.~Belady, \emph{A study of replacement algorithms for a virtual-storage computer}, IBM Systems Journal 5 (1966), 78--101; D.~D.~Sleator and R.~E.~Tarjan, \emph{Amortized efficiency of list update and paging rules}, Communications of the ACM 28 (1985), 202--208; A.~Fiat, R.~M.~Karp, M.~Luby, L.~A.~McGeoch, D.~D.~Sleator and N.~E.~Young, \emph{Competitive paging algorithms}, Journal of Algorithms 12 (1991), 685--699; A.~R.~Karlin, M.~S.~Manasse, L.~A.~McGeoch and S.~Owicki, \emph{Competitive randomized algorithms for nonuniform problems}, Algorithmica 11 (1994), 542--571.
\item N.~Littlestone and M.~K.~Warmuth, \emph{The weighted majority algorithm}, Information and Computation 108 (1994), 212--261; Y.~Freund and R.~E.~Schapire, \emph{A decision-theoretic generalization of on-line learning and an application to boosting}, Journal of Computer and System Sciences 55 (1997), 119--139; S.~Arora, E.~Hazan and S.~Kale, \emph{The multiplicative weights update method: a meta-algorithm and applications}, Theory of Computing 8 (2012), 121--164.
\item D.~P.~Williamson and D.~B.~Shmoys, \emph{The Design of Approximation Algorithms}, Cambridge University Press, 2011, freely available from its authors, is the reference for the first half of the chapter; A.~Borodin and R.~El-Yaniv, \emph{Online Computation and Competitive Analysis}, Cambridge University Press, 1998, for the second.
\end{itemize}
